% p015_b 的电脑重画（通电导线间作用力草图：左为十字坐标线与四个 ×、“↑转”；右为横向电流、竖直通电导线、受力 F大/F小、转）
\begin{tikzpicture}[line width=0.7pt, >={Stealth[length=2.2mm]}]
  % ===== 左：十字线与 × =====
  \draw[->] (0.63,0.8) -- (0.63,3.6);
  \draw[->] (0,2.1) -- (2.45,2.1);
  \node at (1.7,3.1) {$\times$};
  \node at (2.2,3.1) {$\times$};
  \node at (1.7,1.35) {$\times$};
  \node at (2.2,1.35) {$\times$};
  \draw[->] (3.0,1.75) -- (3.0,2.55);
  \node at (3.8,2.1) {转};
  % ===== 右：电流与竖直导线 =====
  \draw[->, line width=1.4pt] (5.3,3.35) -- (7.7,3.35);
  \fill (5.5,3.95) circle (1pt) (6.35,3.95) circle (1pt) (7.2,3.95) circle (1pt);
  \draw[->, thick] (6.33,0.25) -- (6.33,2.8);
  % 上端：向左的力、向右的短粗箭头、F大
  \draw[->, thick] (6.33,2.38) -- (5.5,2.38);
  \draw[->, line width=1.5pt] (6.33,2.52) -- (6.95,2.52);
  \node at (5.3,2.62) {$\times$};
  \node at (6.0,2.62) {$\times$};
  \node at (7.2,2.7) {$\times$};
  \node[anchor=west] at (7.45,2.55) {$F_{\text{大}}$};
  % 逆时针转动箭头与“转”
  \draw[->] (5.25,2.5) .. controls (4.9,2.3) and (4.8,2.0) .. (4.92,1.63);
  \node at (4.62,1.06) {转};
  % 中间的 ×
  \node at (5.85,1.65) {$\times$};
  \node at (6.8,1.65) {$\times$};
  \node at (6.33,0.9) {$\times$};
  % 下端：向左的力、点、F小
  \draw[->, thick] (6.3,0.48) -- (5.7,0.48);
  \fill (5.4,0.48) circle (1pt);
  \node[anchor=west] at (7.1,0.42) {$F_{\text{小}}$};
\end{tikzpicture}
